\chapter{Tabelas Completas de Resultados por Conjunto de Dados}
\label{app:tabelas}

Este apêndice apresenta a tabela mestra do benchmark em forma
desagregada: a métrica primária de cada um dos 14 modelos em cada um
dos 18 conjuntos de dados, no \emph{hold-out} de teste (20\%, semente
42), conforme o protocolo do Capítulo~\ref{chap:metodologia}. São as
Tabelas~\ref{tab:app-binaria}, \ref{tab:app-multiclasse}
e~\ref{tab:app-regressao} --- uma por tarefa, porque a métrica
primária difere entre tarefas e valores de métricas distintas não são
comparáveis entre si.

Em cada tabela, cada linha é um conjunto de dados e cada coluna, um
modelo, na ordem canônica do benchmark (os três GBDTs,
os modelos profundos, as redes KAN e os dois \emph{foundation
models}). A métrica primária segue a definição do protocolo:
\begin{itemize}
  \item \textbf{classificação binária} (10 conjuntos): AUC-ROC ---
    maior é melhor (Tabela~\ref{tab:app-binaria});
  \item \textbf{classificação multiclasse} (3 conjuntos):
    \emph{log-loss} --- menor é melhor
    (Tabela~\ref{tab:app-multiclasse});
  \item \textbf{regressão} (5 conjuntos): RMSE na escala em que a
    suíte de origem distribui o alvo de cada conjunto --- menor é
    melhor (Tabela~\ref{tab:app-regressao}). As escalas diferem entre
    conjuntos (notas da Tabela~\ref{tab:datasets}: o
    \texttt{california\_\allowbreak housing} tem alvo binarizado, de modo que seu
    RMSE é a raiz do escore de Brier; o \texttt{online\_news} tem alvo
    em logaritmo de compartilhamentos), razão pela qual nenhuma
    agregação de RMSE entre conjuntos é feita em lugar algum do
    estudo --- as comparações são sempre por linha ou por \emph{rank}.
\end{itemize}
O melhor valor de cada linha está em \textbf{negrito}. Os valores têm
quatro casas decimais, com vírgula como separador decimal. As duas
células marcadas ``---'' (\texttt{helena} $\times$
TabPFN/TabFM) são as exclusões estruturais do estudo: o conjunto tem
100 classes, acima do limite arquitetural de ambos os
\emph{foundation models}. O estado completo é, portanto, de 250
células medidas em $14 \times 18 = 252$ possíveis.

Quanto à proveniência, todos os números
provêm de
\texttt{results/\allowbreak aggregated/\allowbreak test\_\allowbreak results.\allowbreak csv}, o agregado versionado
do benchmark; as tabelas em si são geradas automaticamente por
\texttt{scripts/\allowbreak gen\_\allowbreak appendix\_\allowbreak tables.\allowbreak py} a partir desse artefato
(nenhum valor é transcrito à mão). Os resultados por \emph{fold} da
validação cruzada interna, que fundamentam as análises estatísticas
do Capítulo~\ref{chap:benchmark}, estão em \texttt{results/\allowbreak aggregated/\allowbreak fold\_\allowbreak results.\allowbreak csv}
--- arquivo de $\sim$300~MB que permanece no repositório local (fora
do controle de versão remoto por tamanho), reproduzível a partir de
\texttt{results/raw} e das sementes registradas.

Quanto ao layout, com 15 colunas (o conjunto de dados mais
os 14 modelos), as tabelas não cabem na orientação de retrato; optou-se
por páginas rotacionadas (\texttt{sidewaystable}, pacote
\texttt{rotating}) com os nomes dos modelos rotacionados no cabeçalho,
mantendo os 14 modelos lado a lado --- a comparação por linha, que é a
leitura principal, fica íntegra numa única tabela por tarefa.

Completam o apêndice os instrumentos restantes do protocolo
estatístico do Capítulo~\ref{chap:metodologia}, todos gerados do mesmo
modo (nenhum valor transcrito à mão): as métricas
\textbf{secundárias} de cada tarefa, resumidas em \emph{rank} médio
por modelo (Tabelas~\ref{tab:app-secundarias-binary},
\ref{tab:app-secundarias-multiclass}
e~\ref{tab:app-secundarias-regression} --- \emph{ranks} porque
\emph{log-loss}, RMSE e MAE têm escalas incomensuráveis entre
conjuntos; os valores por célula estão no artefato), com uma exceção
aberta por inteiro: a estatística \textbf{KS}, métrica padrão da
indústria de crédito, ganha tabela própria com os valores por conjunto
(Tabela~\ref{tab:app-ks}) --- o KS é comensurável entre conjuntos
binários, e é a leitura que conecta o benchmark à prática de risco; as
matrizes de
\textbf{tamanho de efeito} ($d$ de Cohen pareado,
Tabelas~\ref{tab:app-cohensd-binary}
e~\ref{tab:app-cohensd-regression}); os \textbf{intervalos de
confiança de \emph{bootstrap}} da média de AUC na binária
(Tabela~\ref{tab:app-bootstrap-auc}; nas demais tarefas, $N{=}3$ e
$N{=}5$ conjuntos e escalas incomensuráveis tornariam a média entre
conjuntos não interpretável, e o instrumento não é reportado --- a
mesma razão pela qual o estudo não agrega RMSE); os \textbf{custos
consolidados} por modelo (Tabela~\ref{tab:app-custos}); e a
\textbf{variabilidade entre \emph{folds}} da métrica primária
(Tabela~\ref{tab:app-foldvar}), que sustenta a leitura de pior-caso do
Capítulo~\ref{chap:benchmark}.

Por fim, a Tabela~\ref{tab:app-operacionais} abre a matriz de confusão
dos três protagonistas da pergunta central (XGBoost, TabPFN e TabFM) no
\emph{hold-out}: precisão e recall da classe minoritária e acurácia
balanceada, no limiar padrão de 0,5 (argmax), sem nenhuma otimização de
corte --- o retrato operacional que complementa a leitura
\emph{threshold-free} da AUC, conforme o trade-off discutido na
Seção~\ref{sec:fund-metricas-clf} e declarado na
Seção~\ref{sec:met-metricas}. A tabela é computada por
\texttt{scripts/\allowbreak secondary\_\allowbreak metrics\_\allowbreak report.py} sobre as predições por
instância do experimento de complementaridade (Seção~\ref{sec:cap5-ato5}),
com um portão de sanidade próprio: a acurácia recomputada de cada célula
reproduz a do agregado do benchmark (divergência $\leq 10^{-3}$ em todas
as 37 células). F1 e especificidade por célula estão no artefato
(\texttt{secondary\_\allowbreak metrics\_\allowbreak holdout.csv}).

% ============================================================
% ARQUIVO GERADO — NÃO EDITAR À MÃO.
% Gerado por scripts/gen_appendix_tables.py a partir de
% results/aggregated/{test_results,cohens_d_*,bootstrap_ci,
% cv_summary}.csv e results/latency/latency_adult.csv.
% Regenerar: uv run python scripts/gen_appendix_tables.py
% ============================================================


\begin{sidewaystable}[p]
\centering
\small
\setlength{\tabcolsep}{3pt}
\caption[Classificação binária --- resultados completos]{Classificação binária --- AUC-ROC no teste (maior é melhor). Métrica primária por modelo em cada \emph{dataset}; \textbf{negrito} = melhor da linha; ``---'' = célula estruturalmente ausente. Fonte: \texttt{results/\allowbreak aggregated/\allowbreak test\_\allowbreak results.\allowbreak csv}.}
\label{tab:app-binaria}
\begin{tabular}{lcccccccccccccc}
\toprule
\emph{Conjunto} & \rotatebox{90}{XGBoost} & \rotatebox{90}{LightGBM} & \rotatebox{90}{CatBoost} & \rotatebox{90}{MLP} & \rotatebox{90}{RealMLP} & \rotatebox{90}{FT-Transformer} & \rotatebox{90}{SAINT} & \rotatebox{90}{STab} & \rotatebox{90}{TabNet} & \rotatebox{90}{TabM} & \rotatebox{90}{KAN} & \rotatebox{90}{TabKAN} & \rotatebox{90}{TabPFN} & \rotatebox{90}{TabFM} \\
\midrule
\texttt{adult} & 0{,}9301 & 0{,}9299 & 0{,}9297 & 0{,}9128 & 0{,}9190 & 0{,}9165 & 0{,}9201 & 0{,}9172 & 0{,}9121 & 0{,}9161 & 0{,}9151 & 0{,}9140 & 0{,}9204 & \textbf{0{,}9322} \\
\texttt{amazon\_\allowbreak employee} & 0{,}8401 & 0{,}8435 & \textbf{0{,}8869} & 0{,}8208 & 0{,}8068 & 0{,}7835 & 0{,}7793 & 0{,}6769 & 0{,}7888 & 0{,}8108 & 0{,}8171 & 0{,}8125 & 0{,}8317 & 0{,}8604 \\
\texttt{bank\_\allowbreak marketing} & 0{,}9339 & 0{,}9292 & 0{,}9333 & 0{,}9294 & 0{,}9303 & 0{,}9357 & 0{,}9326 & 0{,}9344 & 0{,}9268 & 0{,}9369 & 0{,}9274 & 0{,}9252 & 0{,}9397 & \textbf{0{,}9442} \\
\texttt{cardiovascular\_\allowbreak disease} & 0{,}8002 & 0{,}8001 & 0{,}8001 & 0{,}7981 & 0{,}7980 & 0{,}7993 & 0{,}7998 & 0{,}7991 & 0{,}7805 & 0{,}7991 & 0{,}7979 & 0{,}7955 & 0{,}7997 & \textbf{0{,}8023} \\
\texttt{credit\_g} & 0{,}7667 & 0{,}7656 & 0{,}7720 & 0{,}7744 & 0{,}7362 & \textbf{0{,}7873} & 0{,}7651 & 0{,}7695 & 0{,}7462 & 0{,}7837 & 0{,}7494 & 0{,}7770 & 0{,}7746 & 0{,}7763 \\
\texttt{give\_\allowbreak me\_\allowbreak some\_\allowbreak credit} & 0{,}8640 & 0{,}8640 & 0{,}8646 & 0{,}8212 & 0{,}8625 & 0{,}8270 & 0{,}8273 & 0{,}8274 & 0{,}8175 & 0{,}8268 & 0{,}8190 & 0{,}8187 & 0{,}8638 & \textbf{0{,}8678} \\
\texttt{higgs} & 0{,}8117 & 0{,}8111 & 0{,}8107 & 0{,}8133 & 0{,}8140 & 0{,}8096 & 0{,}8119 & 0{,}8184 & 0{,}8059 & 0{,}8193 & 0{,}8018 & 0{,}7908 & 0{,}8111 & \textbf{0{,}8403} \\
\texttt{magictelescope} & 0{,}9297 & 0{,}9284 & 0{,}9314 & 0{,}9291 & 0{,}9366 & 0{,}9278 & 0{,}9338 & 0{,}9278 & 0{,}9352 & 0{,}9328 & 0{,}9165 & 0{,}9139 & 0{,}9452 & \textbf{0{,}9566} \\
\texttt{phoneme} & 0{,}9429 & 0{,}9416 & 0{,}9454 & 0{,}9213 & 0{,}9482 & 0{,}9362 & 0{,}9358 & 0{,}9138 & 0{,}9232 & 0{,}9388 & 0{,}8952 & 0{,}8939 & 0{,}9529 & \textbf{0{,}9717} \\
\texttt{telco\_\allowbreak customer\_\allowbreak churn} & 0{,}8465 & 0{,}8446 & 0{,}8431 & 0{,}8391 & 0{,}8417 & 0{,}8465 & 0{,}8460 & 0{,}8413 & 0{,}8314 & 0{,}8425 & 0{,}8456 & 0{,}8421 & \textbf{0{,}8490} & 0{,}8233 \\
\bottomrule
\end{tabular}
\end{sidewaystable}

\begin{sidewaystable}[p]
\centering
\small
\setlength{\tabcolsep}{3pt}
\caption[Classificação multiclasse --- resultados completos]{Classificação multiclasse --- log-loss no teste (menor é melhor). Métrica primária por modelo em cada \emph{dataset}; \textbf{negrito} = melhor da linha; ``---'' = célula estruturalmente ausente. Fonte: \texttt{results/\allowbreak aggregated/\allowbreak test\_\allowbreak results.\allowbreak csv}.}
\label{tab:app-multiclasse}
\begin{tabular}{lcccccccccccccc}
\toprule
\emph{Conjunto} & \rotatebox{90}{XGBoost} & \rotatebox{90}{LightGBM} & \rotatebox{90}{CatBoost} & \rotatebox{90}{MLP} & \rotatebox{90}{RealMLP} & \rotatebox{90}{FT-Transformer} & \rotatebox{90}{SAINT} & \rotatebox{90}{STab} & \rotatebox{90}{TabNet} & \rotatebox{90}{TabM} & \rotatebox{90}{KAN} & \rotatebox{90}{TabKAN} & \rotatebox{90}{TabPFN} & \rotatebox{90}{TabFM} \\
\midrule
\texttt{covertype} & 0{,}1934 & 0{,}2009 & 0{,}2035 & 0{,}2118 & 0{,}1894 & 0{,}2076 & 0{,}2289 & 0{,}1786 & 0{,}2423 & 0{,}2067 & 0{,}2496 & 0{,}2577 & 0{,}1644 & \textbf{0{,}0966} \\
\texttt{helena} & 2{,}6010 & 2{,}6265 & 2{,}5800 & 2{,}5260 & 2{,}5144 & 2{,}5420 & 2{,}6090 & 2{,}5342 & 2{,}6758 & \textbf{2{,}5097} & 2{,}6376 & 2{,}7187 & --- & --- \\
\texttt{jannis} & 0{,}6801 & 0{,}6836 & 0{,}6836 & 0{,}6874 & 0{,}6817 & 0{,}6774 & 0{,}6887 & 0{,}6698 & 0{,}7027 & 0{,}6572 & 0{,}7080 & 0{,}7257 & 0{,}6583 & \textbf{0{,}5709} \\
\bottomrule
\end{tabular}
\end{sidewaystable}

\begin{sidewaystable}[p]
\centering
\small
\setlength{\tabcolsep}{3pt}
\caption[Regressão --- resultados completos]{Regressão --- RMSE no teste (menor é melhor). Métrica primária por modelo em cada \emph{dataset}; \textbf{negrito} = melhor da linha; ``---'' = célula estruturalmente ausente. Fonte: \texttt{results/\allowbreak aggregated/\allowbreak test\_\allowbreak results.\allowbreak csv}.}
\label{tab:app-regressao}
\begin{tabular}{lcccccccccccccc}
\toprule
\emph{Conjunto} & \rotatebox{90}{XGBoost} & \rotatebox{90}{LightGBM} & \rotatebox{90}{CatBoost} & \rotatebox{90}{MLP} & \rotatebox{90}{RealMLP} & \rotatebox{90}{FT-Transformer} & \rotatebox{90}{SAINT} & \rotatebox{90}{STab} & \rotatebox{90}{TabNet} & \rotatebox{90}{TabM} & \rotatebox{90}{KAN} & \rotatebox{90}{TabKAN} & \rotatebox{90}{TabPFN} & \rotatebox{90}{TabFM} \\
\midrule
\texttt{california\_\allowbreak housing} & 0{,}2637 & 0{,}2641 & 0{,}2653 & 0{,}3002 & 0{,}2801 & 0{,}2851 & 0{,}2873 & 0{,}2905 & 0{,}3326 & 0{,}2798 & 0{,}3371 & 0{,}3324 & 0{,}2570 & \textbf{0{,}2346} \\
\texttt{kin8nm} & 0{,}1117 & 0{,}1075 & 0{,}0859 & 0{,}0821 & \textbf{0{,}0630} & 0{,}0677 & 0{,}0653 & 0{,}0820 & 0{,}0808 & 0{,}0656 & 0{,}1513 & 0{,}1526 & 0{,}0667 & 0{,}0652 \\
\texttt{online\_news} & 0{,}8308 & 0{,}8320 & 0{,}8331 & 0{,}8523 & 0{,}8401 & 0{,}8432 & 0{,}8464 & 0{,}8415 & 0{,}8594 & 0{,}8482 & 0{,}8457 & 0{,}8460 & 0{,}8284 & \textbf{0{,}8269} \\
\texttt{wine\_quality} & 0{,}5975 & 0{,}6131 & 0{,}6005 & 0{,}6629 & 0{,}6728 & 0{,}6774 & 0{,}6869 & 0{,}7077 & 0{,}6993 & 0{,}6607 & 0{,}6597 & 0{,}6941 & 0{,}6129 & \textbf{0{,}5747} \\
\texttt{year\_\allowbreak prediction} & 8{,}6801 & 8{,}6727 & 8{,}6949 & 9{,}3563 & 8{,}6269 & 8{,}8727 & 8{,}8207 & 8{,}8357 & 8{,}8622 & 8{,}9224 & 8{,}7864 & 8{,}8662 & 8{,}5952 & \textbf{8{,}3634} \\
\bottomrule
\end{tabular}
\end{sidewaystable}

\begin{table}[htbp]
\centering
\small
\caption[Métricas secundárias --- tarefa binária]{Rank médio (1 = melhor) dos 14 modelos em cada métrica da tarefa binária; a primária aparece como referência. Valores por célula em \texttt{results/\allowbreak aggregated/\allowbreak test\_\allowbreak results.\allowbreak csv}. Fonte: idem.}
\label{tab:app-secundarias-binary}
\begin{tabular}{lccccc}
\toprule
\emph{Modelo} & AUC-ROC & Acurácia & F1 & KS & log-loss \\
\midrule
XGBoost & 5{,}00 & 6{,}65 & 6{,}10 & 5{,}00 & 4{,}60 \\
LightGBM & 6{,}30 & 5{,}80 & 7{,}00 & 5{,}65 & 7{,}50 \\
CatBoost & 5{,}30 & 4{,}65 & 5{,}30 & 5{,}50 & 5{,}00 \\
MLP & 9{,}30 & 10{,}05 & 8{,}90 & 9{,}70 & 8{,}80 \\
RealMLP & 7{,}70 & 7{,}50 & 6{,}30 & 7{,}20 & 9{,}80 \\
FT-Transformer & 7{,}40 & 8{,}50 & 7{,}60 & 7{,}00 & 8{,}70 \\
SAINT & 7{,}50 & 7{,}90 & 8{,}20 & 7{,}00 & 6{,}80 \\
STab & 8{,}80 & 8{,}00 & 9{,}20 & 8{,}70 & 7{,}80 \\
TabNet & 11{,}80 & 10{,}30 & 10{,}50 & 11{,}70 & 11{,}20 \\
TabM & 6{,}60 & 7{,}05 & 6{,}90 & 7{,}25 & 5{,}70 \\
KAN & 11{,}00 & 9{,}75 & 10{,}00 & 11{,}00 & 10{,}70 \\
TabKAN & 11{,}40 & 11{,}70 & 11{,}80 & 11{,}50 & 11{,}40 \\
TabPFN & 4{,}20 & 4{,}05 & 5{,}30 & 4{,}80 & 4{,}10 \\
TabFM & \textbf{2{,}70} & \textbf{3{,}10} & \textbf{1{,}90} & \textbf{3{,}00} & \textbf{2{,}90} \\
\bottomrule
\end{tabular}
\end{table}

\begin{table}[htbp]
\centering
\small
\caption[Métricas secundárias --- tarefa multiclasse]{Rank médio (1 = melhor) dos 14 modelos em cada métrica da tarefa multiclasse; a primária aparece como referência. Valores por célula em \texttt{results/\allowbreak aggregated/\allowbreak test\_\allowbreak results.\allowbreak csv}. Fonte: idem.}
\label{tab:app-secundarias-multiclass}
\begin{tabular}{lccccc}
\toprule
\emph{Modelo} & log-loss & Acurácia & F1-macro & F1-ponderado & AUC-OvR \\
\midrule
XGBoost & 6{,}00 & 6{,}33 & 7{,}00 & 5{,}33 & 6{,}67 \\
LightGBM & 8{,}00 & 7{,}67 & 10{,}33 & 7{,}33 & 8{,}33 \\
CatBoost & 7{,}00 & 6{,}67 & 6{,}67 & 6{,}67 & 7{,}33 \\
MLP & 7{,}67 & 8{,}00 & 7{,}67 & 8{,}67 & 8{,}00 \\
RealMLP & 4{,}33 & 3{,}67 & 3{,}00 & 3{,}33 & 4{,}33 \\
FT-Transformer & 6{,}33 & 5{,}33 & 4{,}00 & 4{,}67 & 6{,}00 \\
SAINT & 10{,}00 & 10{,}00 & 9{,}33 & 10{,}00 & 9{,}00 \\
STab & 3{,}67 & 4{,}67 & 6{,}67 & 5{,}67 & 3{,}67 \\
TabNet & 11{,}67 & 11{,}67 & 11{,}67 & 12{,}00 & 11{,}33 \\
TabM & 3{,}67 & 4{,}67 & 5{,}00 & 5{,}67 & 4{,}33 \\
KAN & 12{,}00 & 12{,}33 & 9{,}67 & 12{,}00 & 12{,}00 \\
TabKAN & 13{,}33 & 13{,}00 & 12{,}67 & 12{,}67 & 13{,}00 \\
TabPFN & 2{,}50 & 2{,}00 & 2{,}50 & 2{,}00 & 2{,}00 \\
TabFM & \textbf{1{,}00} & \textbf{1{,}00} & \textbf{1{,}00} & \textbf{1{,}00} & \textbf{1{,}00} \\
\bottomrule
\end{tabular}
\end{table}

\begin{table}[htbp]
\centering
\small
\caption[Métricas secundárias --- tarefa regressão]{Rank médio (1 = melhor) dos 14 modelos em cada métrica da tarefa regressão; a primária aparece como referência. Valores por célula em \texttt{results/\allowbreak aggregated/\allowbreak test\_\allowbreak results.\allowbreak csv}. Fonte: idem.}
\label{tab:app-secundarias-regression}
\begin{tabular}{lccc}
\toprule
\emph{Modelo} & RMSE & MAE & $R^{2}$ \\
\midrule
XGBoost & 5{,}00 & 6{,}40 & 5{,}00 \\
LightGBM & 5{,}60 & 6{,}20 & 5{,}60 \\
CatBoost & 5{,}80 & 6{,}80 & 5{,}80 \\
MLP & 11{,}00 & 9{,}60 & 11{,}00 \\
RealMLP & 5{,}20 & 4{,}40 & 5{,}20 \\
FT-Transformer & 8{,}80 & 8{,}00 & 8{,}80 \\
SAINT & 8{,}40 & 9{,}40 & 8{,}40 \\
STab & 9{,}60 & 10{,}00 & 9{,}60 \\
TabNet & 11{,}40 & 10{,}20 & 11{,}40 \\
TabM & 8{,}40 & 7{,}60 & 8{,}40 \\
KAN & 9{,}80 & 9{,}60 & 9{,}80 \\
TabKAN & 11{,}80 & 12{,}20 & 11{,}80 \\
TabPFN & 3{,}00 & 3{,}20 & 3{,}00 \\
TabFM & \textbf{1{,}20} & \textbf{1{,}40} & \textbf{1{,}20} \\
\bottomrule
\end{tabular}
\end{table}

\begin{sidewaystable}[htbp]
\centering
\small
\setlength{\tabcolsep}{3pt}
\caption[Estatística KS por conjunto --- binária]{Estatística de Kolmogorov--Smirnov (KS; maior é melhor) por modelo em cada conjunto binário --- a métrica de separação padrão da indústria de crédito (Seção~\ref{sec:met-metricas}); \textbf{negrito} = melhor da linha. Fonte: \texttt{results/\allowbreak aggregated/\allowbreak test\_\allowbreak results.\allowbreak csv}.}
\label{tab:app-ks}
\begin{tabular}{lcccccccccccccc}
\toprule
\emph{Conjunto} & \rotatebox{90}{XGBoost} & \rotatebox{90}{LightGBM} & \rotatebox{90}{CatBoost} & \rotatebox{90}{MLP} & \rotatebox{90}{RealMLP} & \rotatebox{90}{FT-Transformer} & \rotatebox{90}{SAINT} & \rotatebox{90}{STab} & \rotatebox{90}{TabNet} & \rotatebox{90}{TabM} & \rotatebox{90}{KAN} & \rotatebox{90}{TabKAN} & \rotatebox{90}{TabPFN} & \rotatebox{90}{TabFM} \\
\midrule
\texttt{adult} & 0{,}6965 & 0{,}6920 & 0{,}6924 & 0{,}6596 & 0{,}6707 & 0{,}6618 & 0{,}6748 & 0{,}6690 & 0{,}6577 & 0{,}6651 & 0{,}6600 & 0{,}6588 & 0{,}6695 & \textbf{0{,}6965} \\
\texttt{amazon\_\allowbreak employee} & 0{,}5376 & 0{,}5674 & \textbf{0{,}6408} & 0{,}5110 & 0{,}5013 & 0{,}4495 & 0{,}4251 & 0{,}2628 & 0{,}4607 & 0{,}5025 & 0{,}4998 & 0{,}4722 & 0{,}5440 & 0{,}6136 \\
\texttt{bank\_\allowbreak marketing} & 0{,}7398 & 0{,}7387 & 0{,}7350 & 0{,}7281 & 0{,}7358 & 0{,}7465 & 0{,}7437 & 0{,}7409 & 0{,}7238 & 0{,}7572 & 0{,}7196 & 0{,}7118 & 0{,}7518 & \textbf{0{,}7671} \\
\texttt{cardiovascular\_\allowbreak disease} & 0{,}4701 & 0{,}4731 & 0{,}4702 & 0{,}4678 & 0{,}4708 & \textbf{0{,}4738} & 0{,}4712 & 0{,}4652 & 0{,}4492 & 0{,}4643 & 0{,}4657 & 0{,}4599 & 0{,}4704 & 0{,}4734 \\
\texttt{credit\_g} & \textbf{0{,}5286} & 0{,}4905 & 0{,}4833 & 0{,}4381 & 0{,}3738 & 0{,}5095 & 0{,}4119 & 0{,}4524 & 0{,}3762 & 0{,}4690 & 0{,}3976 & 0{,}4429 & 0{,}4286 & 0{,}4500 \\
\texttt{give\_\allowbreak me\_\allowbreak some\_\allowbreak credit} & 0{,}5684 & 0{,}5648 & 0{,}5659 & 0{,}5035 & 0{,}5596 & 0{,}5106 & 0{,}5115 & 0{,}5092 & 0{,}4961 & 0{,}5094 & 0{,}4951 & 0{,}4988 & 0{,}5647 & \textbf{0{,}5732} \\
\texttt{higgs} & 0{,}4611 & 0{,}4603 & 0{,}4610 & 0{,}4718 & 0{,}4712 & 0{,}4651 & 0{,}4623 & 0{,}4721 & 0{,}4571 & 0{,}4788 & 0{,}4491 & 0{,}4326 & 0{,}4633 & \textbf{0{,}5128} \\
\texttt{magictelescope} & 0{,}7123 & 0{,}7138 & 0{,}7212 & 0{,}7003 & 0{,}7235 & 0{,}7018 & 0{,}7265 & 0{,}7063 & 0{,}7265 & 0{,}7138 & 0{,}6831 & 0{,}6786 & 0{,}7519 & \textbf{0{,}7930} \\
\texttt{phoneme} & 0{,}7703 & 0{,}7805 & 0{,}7692 & 0{,}7087 & 0{,}7599 & 0{,}7503 & 0{,}7392 & 0{,}6882 & 0{,}7179 & 0{,}7458 & 0{,}6435 & 0{,}6548 & 0{,}7834 & \textbf{0{,}8270} \\
\texttt{telco\_\allowbreak customer\_\allowbreak churn} & 0{,}5381 & 0{,}5285 & 0{,}5340 & 0{,}5171 & 0{,}5300 & 0{,}5288 & 0{,}5347 & 0{,}5301 & 0{,}5079 & 0{,}5262 & \textbf{0{,}5415} & 0{,}5351 & 0{,}5389 & 0{,}5150 \\
\bottomrule
\end{tabular}
\end{sidewaystable}

\begin{sidewaystable}[p]
\centering
\scriptsize
\setlength{\tabcolsep}{3pt}
\caption[Tamanhos de efeito ($d$ de Cohen) --- binária]{$d$ de Cohen pareado (linha $-$ coluna) na tarefa binária (AUC-ROC); positivo = linha melhor (AUC maior). Fonte: \texttt{results/\allowbreak aggregated/\allowbreak cohens\_\allowbreak d\_\allowbreak binary.\allowbreak csv}.}
\label{tab:app-cohensd-binary}
\begin{tabular}{lcccccccccccccc}
\toprule
 & \rotatebox{90}{XGBoost} & \rotatebox{90}{LightGBM} & \rotatebox{90}{CatBoost} & \rotatebox{90}{MLP} & \rotatebox{90}{RealMLP} & \rotatebox{90}{FT-Transformer} & \rotatebox{90}{SAINT} & \rotatebox{90}{STab} & \rotatebox{90}{TabNet} & \rotatebox{90}{TabM} & \rotatebox{90}{KAN} & \rotatebox{90}{TabKAN} & \rotatebox{90}{TabPFN} & \rotatebox{90}{TabFM} \\
\midrule
XGBoost & --- & $0{,}38$ & $-0{,}35$ & $0{,}72$ & $0{,}52$ & $0{,}44$ & $0{,}55$ & $0{,}47$ & $1{,}14$ & $0{,}35$ & $1{,}10$ & $0{,}99$ & $-0{,}28$ & $-0{,}68$ \\
LightGBM & $-0{,}38$ & --- & $-0{,}44$ & $0{,}63$ & $0{,}43$ & $0{,}38$ & $0{,}48$ & $0{,}44$ & $1{,}01$ & $0{,}29$ & $1{,}02$ & $0{,}91$ & $-0{,}33$ & $-0{,}73$ \\
CatBoost & $0{,}35$ & $0{,}44$ & --- & $0{,}69$ & $0{,}47$ & $0{,}43$ & $0{,}48$ & $0{,}45$ & $0{,}85$ & $0{,}41$ & $0{,}97$ & $0{,}91$ & $0{,}15$ & $-0{,}31$ \\
MLP & $-0{,}72$ & $-0{,}63$ & $-0{,}69$ & --- & $-0{,}16$ & $-0{,}07$ & $0{,}05$ & $0{,}29$ & $0{,}72$ & $-0{,}68$ & $0{,}67$ & $0{,}71$ & $-0{,}91$ & $-1{,}02$ \\
RealMLP & $-0{,}52$ & $-0{,}43$ & $-0{,}47$ & $0{,}16$ & --- & $0{,}11$ & $0{,}23$ & $0{,}38$ & $0{,}84$ & $-0{,}07$ & $0{,}49$ & $0{,}41$ & $-0{,}75$ & $-0{,}91$ \\
FT-Transformer & $-0{,}44$ & $-0{,}38$ & $-0{,}43$ & $0{,}07$ & $-0{,}11$ & --- & $0{,}23$ & $0{,}43$ & $0{,}74$ & $-0{,}41$ & $0{,}41$ & $0{,}49$ & $-0{,}65$ & $-0{,}72$ \\
SAINT & $-0{,}55$ & $-0{,}48$ & $-0{,}48$ & $-0{,}05$ & $-0{,}23$ & $-0{,}23$ & --- & $0{,}39$ & $0{,}94$ & $-0{,}49$ & $0{,}34$ & $0{,}34$ & $-0{,}77$ & $-0{,}81$ \\
STab & $-0{,}47$ & $-0{,}44$ & $-0{,}45$ & $-0{,}29$ & $-0{,}38$ & $-0{,}43$ & $-0{,}39$ & --- & $-0{,}11$ & $-0{,}44$ & $-0{,}12$ & $-0{,}12$ & $-0{,}55$ & $-0{,}62$ \\
TabNet & $-1{,}14$ & $-1{,}01$ & $-0{,}85$ & $-0{,}72$ & $-0{,}84$ & $-0{,}74$ & $-0{,}94$ & $0{,}11$ & --- & $-1{,}28$ & $-0{,}11$ & $-0{,}08$ & $-1{,}54$ & $-1{,}40$ \\
TabM & $-0{,}35$ & $-0{,}29$ & $-0{,}41$ & $0{,}68$ & $0{,}07$ & $0{,}41$ & $0{,}49$ & $0{,}44$ & $1{,}28$ & --- & $0{,}75$ & $0{,}84$ & $-0{,}59$ & $-0{,}78$ \\
KAN & $-1{,}10$ & $-1{,}02$ & $-0{,}97$ & $-0{,}67$ & $-0{,}49$ & $-0{,}41$ & $-0{,}34$ & $0{,}12$ & $0{,}11$ & $-0{,}75$ & --- & $0{,}02$ & $-1{,}08$ & $-1{,}07$ \\
TabKAN & $-0{,}99$ & $-0{,}91$ & $-0{,}91$ & $-0{,}71$ & $-0{,}41$ & $-0{,}49$ & $-0{,}34$ & $0{,}12$ & $0{,}08$ & $-0{,}84$ & $-0{,}02$ & --- & $-1{,}05$ & $-1{,}00$ \\
TabPFN & $0{,}28$ & $0{,}33$ & $-0{,}15$ & $0{,}91$ & $0{,}75$ & $0{,}65$ & $0{,}77$ & $0{,}55$ & $1{,}54$ & $0{,}59$ & $1{,}08$ & $1{,}05$ & --- & $-0{,}55$ \\
TabFM & $0{,}68$ & $0{,}73$ & $0{,}31$ & $1{,}02$ & $0{,}91$ & $0{,}72$ & $0{,}81$ & $0{,}62$ & $1{,}40$ & $0{,}78$ & $1{,}07$ & $1{,}00$ & $0{,}55$ & --- \\
\bottomrule
\end{tabular}
\end{sidewaystable}

\begin{sidewaystable}[p]
\centering
\scriptsize
\setlength{\tabcolsep}{3pt}
\caption[Tamanhos de efeito ($d$ de Cohen) --- regressão]{$d$ de Cohen pareado (linha $-$ coluna) na tarefa regressão (RMSE); positivo = linha \emph{pior} (RMSE maior). Fonte: \texttt{results/\allowbreak aggregated/\allowbreak cohens\_\allowbreak d\_\allowbreak regression.\allowbreak csv}.}
\label{tab:app-cohensd-regression}
\begin{tabular}{lcccccccccccccc}
\toprule
 & \rotatebox{90}{XGBoost} & \rotatebox{90}{LightGBM} & \rotatebox{90}{CatBoost} & \rotatebox{90}{MLP} & \rotatebox{90}{RealMLP} & \rotatebox{90}{FT-Transformer} & \rotatebox{90}{SAINT} & \rotatebox{90}{STab} & \rotatebox{90}{TabNet} & \rotatebox{90}{TabM} & \rotatebox{90}{KAN} & \rotatebox{90}{TabKAN} & \rotatebox{90}{TabPFN} & \rotatebox{90}{TabFM} \\
\midrule
XGBoost & --- & $-0{,}13$ & $0{,}05$ & $-0{,}52$ & $0{,}00$ & $-0{,}58$ & $-0{,}62$ & $-0{,}72$ & $-0{,}88$ & $-0{,}53$ & $-1{,}71$ & $-1{,}24$ & $0{,}61$ & $0{,}64$ \\
LightGBM & $0{,}13$ & --- & $0{,}12$ & $-0{,}51$ & $0{,}03$ & $-0{,}56$ & $-0{,}61$ & $-0{,}71$ & $-0{,}86$ & $-0{,}51$ & $-1{,}55$ & $-1{,}18$ & $0{,}78$ & $0{,}67$ \\
CatBoost & $-0{,}05$ & $-0{,}12$ & --- & $-0{,}54$ & $-0{,}01$ & $-0{,}69$ & $-0{,}77$ & $-0{,}86$ & $-1{,}06$ & $-0{,}60$ & $-2{,}06$ & $-1{,}43$ & $0{,}54$ & $0{,}60$ \\
MLP & $0{,}52$ & $0{,}51$ & $0{,}54$ & --- & $0{,}48$ & $0{,}47$ & $0{,}46$ & $0{,}42$ & $0{,}36$ & $0{,}50$ & $0{,}35$ & $0{,}31$ & $0{,}55$ & $0{,}56$ \\
RealMLP & $-0{,}00$ & $-0{,}03$ & $0{,}01$ & $-0{,}48$ & --- & $-0{,}49$ & $-0{,}54$ & $-0{,}63$ & $-0{,}75$ & $-0{,}44$ & $-0{,}86$ & $-0{,}87$ & $1{,}03$ & $0{,}78$ \\
FT-Transformer & $0{,}58$ & $0{,}56$ & $0{,}69$ & $-0{,}47$ & $0{,}49$ & --- & $0{,}31$ & $-0{,}09$ & $-0{,}85$ & $-0{,}24$ & $-0{,}10$ & $-0{,}78$ & $0{,}67$ & $0{,}64$ \\
SAINT & $0{,}62$ & $0{,}61$ & $0{,}77$ & $-0{,}46$ & $0{,}54$ & $-0{,}31$ & --- & $-0{,}95$ & $-1{,}56$ & $-0{,}28$ & $-0{,}28$ & $-1{,}05$ & $0{,}76$ & $0{,}68$ \\
STab & $0{,}72$ & $0{,}71$ & $0{,}86$ & $-0{,}42$ & $0{,}63$ & $0{,}09$ & $0{,}95$ & --- & $-0{,}75$ & $-0{,}08$ & $-0{,}08$ & $-0{,}82$ & $0{,}83$ & $0{,}72$ \\
TabNet & $0{,}88$ & $0{,}86$ & $1{,}06$ & $-0{,}36$ & $0{,}75$ & $0{,}85$ & $1{,}56$ & $0{,}75$ & --- & $0{,}26$ & $0{,}20$ & $-0{,}33$ & $0{,}94$ & $0{,}78$ \\
TabM & $0{,}53$ & $0{,}51$ & $0{,}60$ & $-0{,}50$ & $0{,}44$ & $0{,}24$ & $0{,}28$ & $0{,}08$ & $-0{,}26$ & --- & $-0{,}01$ & $-0{,}42$ & $0{,}61$ & $0{,}61$ \\
KAN & $1{,}71$ & $1{,}55$ & $2{,}06$ & $-0{,}35$ & $0{,}86$ & $0{,}10$ & $0{,}28$ & $0{,}08$ & $-0{,}20$ & $0{,}01$ & --- & $-0{,}62$ & $1{,}28$ & $0{,}90$ \\
TabKAN & $1{,}24$ & $1{,}18$ & $1{,}43$ & $-0{,}31$ & $0{,}87$ & $0{,}78$ & $1{,}05$ & $0{,}82$ & $0{,}33$ & $0{,}42$ & $0{,}62$ & --- & $1{,}10$ & $0{,}86$ \\
TabPFN & $-0{,}61$ & $-0{,}78$ & $-0{,}54$ & $-0{,}55$ & $-1{,}03$ & $-0{,}67$ & $-0{,}76$ & $-0{,}83$ & $-0{,}94$ & $-0{,}61$ & $-1{,}28$ & $-1{,}10$ & --- & $0{,}60$ \\
TabFM & $-0{,}64$ & $-0{,}67$ & $-0{,}60$ & $-0{,}56$ & $-0{,}78$ & $-0{,}64$ & $-0{,}68$ & $-0{,}72$ & $-0{,}78$ & $-0{,}61$ & $-0{,}90$ & $-0{,}86$ & $-0{,}60$ & --- \\
\bottomrule
\end{tabular}
\end{sidewaystable}

\begin{table}[htbp]
\centering
\small
\caption[IC bootstrap da média de AUC --- binária]{Média de AUC-ROC na tarefa binária com IC95\% de \emph{bootstrap} percentil ($B{=}10\,000$) sobre os conjuntos. Fonte: \texttt{results/\allowbreak aggregated/\allowbreak bootstrap\_\allowbreak ci.\allowbreak csv}.}
\label{tab:app-bootstrap-auc}
\begin{tabular}{lcc}
\toprule
\emph{Modelo} & Média & IC95\% \\
\midrule
XGBoost & 0{,}8666 & [0{,}8282; 0{,}9037] \\
LightGBM & 0{,}8658 & [0{,}8279; 0{,}9025] \\
CatBoost & 0{,}8717 & [0{,}8335; 0{,}9077] \\
MLP & 0{,}8560 & [0{,}8202; 0{,}8920] \\
RealMLP & 0{,}8593 & [0{,}8159; 0{,}9008] \\
FT-Transformer & 0{,}8569 & [0{,}8190; 0{,}8954] \\
SAINT & 0{,}8552 & [0{,}8145; 0{,}8958] \\
STab & 0{,}8426 & [0{,}7916; 0{,}8888] \\
TabNet & 0{,}8468 & [0{,}8049; 0{,}8886] \\
TabM & 0{,}8607 & [0{,}8236; 0{,}8983] \\
KAN & 0{,}8485 & [0{,}8119; 0{,}8843] \\
TabKAN & 0{,}8484 & [0{,}8142; 0{,}8822] \\
TabPFN & 0{,}8688 & [0{,}8293; 0{,}9079] \\
TabFM & \textbf{0{,}8775} & [0{,}8365; 0{,}9181] \\
\bottomrule
\end{tabular}
\end{table}

\begin{table}[htbp]
\centering
\small
\caption[Custos consolidados por modelo]{Custos medianos por modelo: treino (s, 18 conjuntos), tuning (s, excl.\ \texttt{adult} --- Seção~\ref{sec:met-custo}) e latência (\textmu s/linha, \texttt{adult}). Fonte: \texttt{test\_\allowbreak results.\allowbreak csv} + \texttt{latency\_\allowbreak adult.\allowbreak csv}.}
\label{tab:app-custos}
\begin{tabular}{lrrr}
\toprule
\emph{Modelo} & Treino (s) & Tuning (s) & Latência (\textmu s/linha) \\
\midrule
XGBoost & 0,69 & 157 & 0,33 \\
LightGBM & 0,66 & 104 & 1,91 \\
CatBoost & 3,76 & 646 & 1,43 \\
MLP & 12,95 & 1.117 & 2,37 \\
RealMLP & 36,90 & 1.352 & 1,52 \\
FT-Transformer & 87,22 & 4.505 & 28,87 \\
SAINT & 202,19 & 9.831 & 44,69 \\
STab & 198,85 & 14.622 & 8.644,82 \\
TabNet & 64,97 & 5.928 & 17,09 \\
TabM & 33,76 & 2.242 & 3,48 \\
KAN & 11,94 & 767 & 3,07 \\
TabKAN & 9,15 & 690 & 2,32 \\
TabPFN & 1,35 & --- & 7.416,10 \\
TabFM & 0,18 & --- & 43.262,04 \\
\bottomrule
\end{tabular}
\end{table}

\begin{table}[htbp]
\centering
\small
\caption[Variabilidade entre \emph{folds}]{Mediana, sobre os conjuntos da tarefa, do coeficiente de variação entre \emph{folds} da métrica primária (\%; menor = mais estável). Fonte: \texttt{results/\allowbreak aggregated/\allowbreak cv\_\allowbreak summary.\allowbreak csv}.}
\label{tab:app-foldvar}
\begin{tabular}{lccc}
\toprule
\emph{Modelo} & Binária (AUC) & Multiclasse (log-loss) & Regressão (RMSE) \\
\midrule
XGBoost & 0{,}49 & 0{,}79 & 1{,}96 \\
LightGBM & 0{,}52 & 0{,}73 & 1{,}79 \\
CatBoost & 0{,}35 & 0{,}85 & 2{,}11 \\
MLP & 0{,}44 & 0{,}65 & 3{,}75 \\
RealMLP & 0{,}51 & 1{,}22 & 1{,}61 \\
FT-Transformer & 0{,}46 & 0{,}77 & 1{,}46 \\
SAINT & 0{,}62 & 1{,}11 & 1{,}61 \\
STab & 0{,}46 & 1{,}32 & 3{,}60 \\
TabNet & 0{,}73 & 0{,}98 & 2{,}47 \\
TabM & 0{,}48 & 0{,}75 & 1{,}86 \\
KAN & 0{,}43 & 1{,}10 & 1{,}87 \\
TabKAN & 0{,}50 & 0{,}80 & 1{,}69 \\
TabPFN & 0{,}37 & 1{,}60 & 1{,}48 \\
TabFM & 0{,}36 & 1{,}57 & 1{,}46 \\
\bottomrule
\end{tabular}
\end{table}

% ============================================================
% ARQUIVO GERADO - NAO EDITAR A MAO.
% Gerado por scripts/secondary_metrics_report.py a partir de
% results/aggregated/secondary_metrics_holdout.csv.
% Regenerar: uv run python scripts/secondary_metrics_report.py
% ============================================================

\begin{table}[htbp]
\centering
\footnotesize
\setlength{\tabcolsep}{3.5pt}
\caption[Métricas operacionais no \emph{hold-out} --- binárias]{Métricas operacionais no \emph{hold-out} (limiar 0,5; classe de interesse = minoritária). Fonte: \texttt{secondary\_metrics\_holdout.csv}.}
\label{tab:app-operacionais}
\begin{tabular}{lr ccc ccc ccc}
\toprule
 & & \multicolumn{3}{c}{Precisão (min.)} & \multicolumn{3}{c}{Recall (min.)} & \multicolumn{3}{c}{Acurácia balanceada} \\
\cmidrule(lr){3-5}\cmidrule(lr){6-8}\cmidrule(lr){9-11}
\emph{Conjunto} & \%min. & XGB & PFN & FM & XGB & PFN & FM & XGB & PFN & FM \\
\midrule
\texttt{adult} & 23{,}9 & 0{,}795 & 0{,}783 & \textbf{0{,}803} & 0{,}648 & 0{,}618 & \textbf{0{,}651} & 0{,}798 & 0{,}782 & \textbf{0{,}800} \\
\texttt{amazon\_\allowbreak employee} & 5{,}8 & \textbf{0{,}706} & 0{,}659 & 0{,}621 & 0{,}032 & 0{,}214 & \textbf{0{,}298} & 0{,}515 & 0{,}603 & \textbf{0{,}643} \\
\texttt{bank\_\allowbreak marketing} & 11{,}7 & 0{,}647 & \textbf{0{,}676} & 0{,}660 & 0{,}466 & 0{,}526 & \textbf{0{,}568} & 0{,}716 & 0{,}747 & \textbf{0{,}765} \\
\texttt{cardiovascular\_\allowbreak disease} & 50{,}0 & \textbf{0{,}754} & 0{,}751 & \textbf{0{,}754} & 0{,}693 & \textbf{0{,}700} & 0{,}698 & 0{,}734 & 0{,}734 & \textbf{0{,}735} \\
\texttt{credit\_\allowbreak g} & 30{,}0 & \textbf{0{,}667} & 0{,}659 & 0{,}630 & 0{,}433 & 0{,}450 & \textbf{0{,}483} & 0{,}670 & 0{,}675 & \textbf{0{,}681} \\
\texttt{give\_\allowbreak me\_\allowbreak some\_\allowbreak credit} & 6{,}7 & \textbf{0{,}615} & 0{,}592 & 0{,}596 & 0{,}162 & \textbf{0{,}181} & 0{,}163 & 0{,}578 & \textbf{0{,}586} & 0{,}578 \\
\texttt{higgs} & 47{,}1 & 0{,}716 & 0{,}715 & \textbf{0{,}746} & 0{,}708 & 0{,}700 & \textbf{0{,}732} & 0{,}729 & 0{,}725 & \textbf{0{,}755} \\
\texttt{magictelescope} & 50{,}0 & 0{,}824 & 0{,}848 & \textbf{0{,}865} & 0{,}886 & 0{,}907 & \textbf{0{,}928} & 0{,}848 & 0{,}873 & \textbf{0{,}892} \\
\texttt{phoneme} & 29{,}3 & 0{,}808 & 0{,}827 & \textbf{0{,}876} & 0{,}785 & 0{,}801 & \textbf{0{,}849} & 0{,}854 & 0{,}866 & \textbf{0{,}899} \\
\texttt{telco\_\allowbreak customer\_\allowbreak churn} & 26{,}5 & 0{,}642 & \textbf{0{,}678} & 0{,}559 & 0{,}570 & 0{,}484 & \textbf{0{,}663} & 0{,}727 & 0{,}700 & \textbf{0{,}737} \\
\midrule
 & & \multicolumn{3}{c}{Precisão macro} & \multicolumn{3}{c}{Recall macro} & \multicolumn{3}{c}{F1 macro} \\
\cmidrule(lr){3-5}\cmidrule(lr){6-8}\cmidrule(lr){9-11}
\texttt{covertype} & --- & 0{,}908 & 0{,}907 & \textbf{0{,}943} & 0{,}857 & 0{,}887 & \textbf{0{,}927} & 0{,}879 & 0{,}896 & \textbf{0{,}935} \\
\texttt{jannis} & --- & 0{,}668 & 0{,}682 & \textbf{0{,}742} & 0{,}552 & 0{,}579 & \textbf{0{,}663} & 0{,}566 & 0{,}599 & \textbf{0{,}690} \\
\bottomrule
\end{tabular}
\end{table}

